% !TeX program = xelatex
% Mathematical companion to the original R report and Simulation2.
% Shared XeLaTeX / KOMA-Script report style for composite_iq_checked_claims.
% Modeled on iq_test_differences_report.tex supplied by Anders Hellström.
% Requires XeLaTeX; loaded before \begin{document} from a report's main file.
% Reports keep their own mathematics, references, assumptions and AI appendix.
\documentclass[11pt,a4paper,oneside,open=any,parskip=half,headings=normal,
  numbers=noenddot,listof=totoc,bibliography=totoc]{scrbook}
\usepackage{iftex}
\RequireXeTeX
\usepackage[margin=18mm,top=19mm,bottom=20mm,headheight=16pt,headsep=7mm]{geometry}
\usepackage{fontspec}
\setmainfont{lmroman10-regular.otf}[BoldFont=lmroman10-bold.otf,
  ItalicFont=lmroman10-italic.otf,BoldItalicFont=lmroman10-bolditalic.otf]
\setsansfont{lmsans10-regular.otf}[BoldFont=lmsans10-bold.otf,
  ItalicFont=lmsans10-oblique.otf,BoldItalicFont=lmsans10-boldoblique.otf]
\setmonofont{DejaVu Sans Mono}[Scale=0.79]
\usepackage{mathtools,amsthm}
\usepackage{unicode-math}
\setmathfont{latinmodern-math.otf}
\usepackage[protrusion=true,expansion=false]{microtype}
\usepackage{xcolor}
\usepackage[most]{tcolorbox}
\usepackage{tabularray}
\UseTblrLibrary{booktabs}
\usepackage{booktabs,longtable,array}
\usepackage{graphicx,float,enumitem}
\usepackage{scrhack,listings}
\usepackage{needspace,placeins}
\usepackage[round,authoryear]{natbib}
\usepackage{xurl}
\usepackage{scrlayer-scrpage}
\usepackage{hyperref}
\usepackage{bookmark}
\usepackage[nameinlink,noabbrev]{cleveref}

% Centralized palette and link appearance from the reference report.
\definecolor{ink}{HTML}{183B56}
\definecolor{teal}{HTML}{007F82}
\definecolor{rust}{HTML}{B75D36}
\definecolor{pale}{HTML}{EFF5F7}
\newcommand{\ProjectAuthor}{Anders Hellström}
\hypersetup{unicode=true,colorlinks=true,linkcolor=ink,citecolor=teal,
  urlcolor=teal,pdfauthor={Anders Hellström; AI-assisted research and drafting}}
\clearpairofpagestyles
\automark[chapter]{chapter}
\renewcommand*{\chaptermarkformat}{}
\ihead{\small\sffamily \headmark}
\ohead{\small\sffamily \headmark}
\cfoot*{\pagemark}
\pagestyle{scrheadings}
\setkomafont{disposition}{\sffamily\bfseries\color{ink}}
\setkomafont{pageheadfoot}{\normalfont\sffamily\small\color{ink}}
\RedeclareSectionCommand[beforeskip=0pt,afterskip=14pt]{chapter}
\setlength{\emergencystretch}{2em}
\setlength{\parfillskip}{0pt plus 1fil}
\setcounter{tocdepth}{2}
\setcounter{secnumdepth}{2}
\setlist{nosep,leftmargin=*}
\setlength{\bibsep}{3pt}
\renewcommand{\bibfont}{\small}

% Preserve names and numbering used by the existing mathematical books.
\DeclareMathOperator{\E}{E}
\DeclareMathOperator{\Var}{Var}
\DeclareMathOperator{\Cov}{Cov}
\DeclareMathOperator{\SD}{SD}
\DeclareMathOperator{\Corr}{Corr}
\newcommand{\T}{\mathsf{T}}
\newcommand{\given}{\,\vert\,}
\newcommand{\one}{\mathbf{1}}
\newtheorem{theorem}{Theorem}[chapter]
\newtheorem{proposition}[theorem]{Proposition}
\newtheorem{lemma}[theorem]{Lemma}
\newtheorem{corollary}[theorem]{Corollary}
\crefname{chapter}{chapter}{chapters}
\Crefname{chapter}{Chapter}{Chapters}
\crefname{proposition}{proposition}{propositions}
\Crefname{proposition}{Proposition}{Propositions}
\crefname{corollary}{corollary}{corollaries}
\Crefname{corollary}{Corollary}{Corollaries}
\counterwithout{equation}{chapter}

% Boxes, code and front matter share the reference report's design.
\newtcolorbox{resultbox}[1][]{colback=pale,colframe=ink,boxrule=.7pt,
  arc=1mm,left=3mm,right=3mm,top=2mm,bottom=2mm,breakable,#1}
\newtcolorbox{disclosurebox}[1][]{colback=rust!7,colframe=rust,boxrule=1.1pt,
  arc=0mm,left=3mm,right=3mm,top=3mm,bottom=3mm,breakable,#1}
\lstset{basicstyle=\ttfamily\footnotesize,breaklines=true,columns=fullflexible,
  keepspaces=true,showstringspaces=false,numbers=left,numberstyle=\tiny\color{gray},
  numbersep=6pt,xleftmargin=15pt,frame=single,rulecolor=\color{ink!25},
  framerule=.3pt,framesep=5pt,backgroundcolor=\color{pale!45},
  keywordstyle=\color{ink}\bfseries,commentstyle=\color{teal},
  stringstyle=\color{rust},tabsize=2,upquote=true}
\newenvironment{abstract}{\addchap{Overview}}{}

% Report-specific summaries are supplied in each book's preamble.
\newcommand{\ProjectReportType}{MATHEMATICAL AND COMPUTATIONAL REPORT}
\newcommand{\ProjectReportHighlights}{\textbf{Scope:} mathematical deductions,
  numerical approximations and empirical claims are distinguished.}
\makeatletter
\newcommand*{\ProjectReportTitlePage}{%
  \begin{titlepage}
    \thispagestyle{empty}
    \vspace*{6mm}
    {\sffamily\color{teal}\large \ProjectReportType\par}
    \vspace{10mm}
    {\sffamily\bfseries\color{ink}\fontsize{27}{32}\selectfont\@title\par}
    \vspace{6mm}
    {\Large\@subtitle\par}
    \vspace{8mm}
    \begin{disclosurebox}[title={\sffamily\bfseries AI-GENERATED REPORT: READ THIS FIRST}]
      This document and substantial supporting mathematical exposition, programs,
      verification and typesetting were prepared with AI assistance at the request
      of Anders Hellström. The reported numerical results are conditional on
      explicit models and evidence; automated checks are not independent human
      peer review or empirical psychometric validation. Full AI-use details
      appear in the final appendix.
    \end{disclosurebox}
    \vspace{5mm}
    \begin{resultbox}[title={\sffamily\bfseries Key findings and interpretation}]
      \ProjectReportHighlights
    \end{resultbox}
    \vspace{3mm}
    \begin{resultbox}[title={\sffamily\bfseries Project repository and reproducibility}]
      \textbf{Repository:}\quad
      \url{https://github.com/AndersH3/composite_iq_checked_claims}

      Source, proof and simulation files are organized in the repository;
      this report specifies its own assumptions, references and verification scope.
    \end{resultbox}
    \vfill
    Prepared at the request of \textbf{\ProjectAuthor}\\
    Report date: \textbf{\@date}\\
    Authoring workflow: OpenAI-assisted research and typesetting\\
    Format: XeLaTeX, KOMA-Script \texttt{scrbook}; A4, narrow margins
  \end{titlepage}%
}
\makeatother

\title{Differences and Ranges of IQ Scores}
\subtitle{Independent individuals, correlated tests, and conditional discrepancies}
\author{\ProjectAuthor}
\date{26 November 2026}
\renewcommand{\ProjectReportHighlights}{For independent normal IQ scores the
mean absolute gap is about 16.93 points; for correlation 0.6 it is 10.70.
The expected five-test range is 22.07 under joint Gaussian equicorrelation.
These are model deductions, not observed psychometric results.}
\begin{document}
\frontmatter
\ProjectReportTitlePage
\begin{abstract}
We distinguish interindividual comparisons from intraindividual test
discrepancies, derive the Gaussian answers using standard probability
results, and expose their dependence on the joint distribution.
A separate simulation book documents numerical checks and their status.
\end{abstract}
\tableofcontents
\mainmatter
\chapter{Model, estimands and assumptions}
The upstream analyses are \citep{iqsource,simsource}. Let $X_j$ denote test
$j$ for one person, or let $X_1,X_2$ denote scores of independently sampled
people on the same test. Define $D=X_1-X_2$, $A=|D|$, and
$R_k=\max_j X_j-\min_j X_j$, with $k\ge2$. The standard example uses
$\mu=100$, $\sigma=15$ and $\rho=0.6$ for distinct tests; independence
between people is a separate sampling hypothesis. A test correlation is
not a reliability coefficient without a measurement model.
All propositions below concern synthetic distributions.
% First two moments require no normality.
\begin{proposition}[Moment identities]\label{prop:moments}
If two scores have finite second moments, common mean and common variance
$\sigma^2>0$, and correlation $\rho$, then
$\E D=0$, $\E D^2=2\sigma^2(1-\rho)$ and
$\E A\le \sigma\sqrt{2(1-\rho)}$.
\end{proposition}
\begin{proof}
Use linearity of expectation, bilinearity of covariance and
Cauchy--Schwarz applied to $|D|$ and the constant one.
\end{proof}
% Projection representation also covers singular and negative correlations.
\begin{proposition}[Admissible equicorrelation]\label{prop:domain}
The matrix $C=(1-\rho)I_k+\rho\one\one^\T$ is positive semidefinite
exactly when $-1/(k-1)\le\rho\le1$.
For iid standard normals $Z_j$, with $\bar Z=k^{-1}\sum_j Z_j$, the vector
\[
X_j=\mu+\sigma\{\sqrt{1-\rho}(Z_j-\bar Z)
+\sqrt{1+(k-1)\rho}\,\bar Z\}
\]
has the Gaussian law with covariance $\sigma^2C$.
\end{proposition}
\begin{proof}
The orthogonal projections $P=\one\one^\T/k$ and $I-P$ have complementary
ranges of dimensions one and $k-1$. On them $C$ acts by
$1+(k-1)\rho$ and $1-\rho$. Both eigenvalues must be nonnegative.
The displayed linear map has product with its transpose equal to
$\sigma^2C$ because cross-products of the projections vanish.
Closure of Gaussian laws under linear maps and uniqueness by their
characteristic functions give the claim \citep{anderson}.
\end{proof}
\chapter{Exact Gaussian differences and ranges}
% Folded-normal expectation and distribution.
\begin{proposition}[Absolute two-score gap]\label{prop:gap}
Under joint normality, let $\tau^2=2\sigma^2(1-\rho)$.
For $\tau>0$ and $a\ge0$,
\[
\E A=\tau\sqrt{2/\pi},\quad
\mathbb{P}(A\le a)=2\Phi(a/\tau)-1,\quad
\Var(A)=\tau^2(1-2/\pi).
\]
For $\tau=0$ the gap vanishes almost surely.
\end{proposition}
\begin{proof}
The contrast is $N(0,\tau^2)$ by \cref{prop:moments}.
Integrating the normal density over $[-a,a]$ gives the CDF; integrating
$2x\phi(x)$ on the positive half-line gives $\sqrt{2/\pi}$.
Since $A^2=D^2$, subtract the squared mean from its second moment.
Zero variance with zero mean forces $D=0$ almost surely.
\end{proof}
Independent people with normal scores correspond to $\rho=0$, giving
$30/\sqrt{\pi}\approx16.9257$; correlated tests with $\rho=0.6$
give $30\sqrt{0.4/\pi}\approx10.7047$. These are not RMS gaps.
% The common projection cancels pointwise, including negative rho.
\begin{proposition}[Range distribution and expectation]\label{prop:range}
Put $s=\sigma\sqrt{1-\rho}$ and let $M_k=\max_j Z_j$. Then
\[
\E R_k=2s\,a_k,\qquad a_k=k\int_{\mathbb R}z\phi(z)\Phi(z)^{k-1}\,dz.
\]
For $s>0$ and $r\ge0$,
\[
\mathbb{P}(R_k\le r)=k\int_{\mathbb R}\phi(z)
[\Phi(z+r/s)-\Phi(z)]^{k-1}\,dz.
\]
\end{proposition}
\begin{proof}
By \cref{prop:domain}, common additive terms cancel and $R_k=s(M_k-\min Z_j)$.
Sign symmetry and integrability (domination by $\sum |Z_j|$) imply the
expectation identity. Independence yields $\mathbb{P}(M_k\le z)=\Phi(z)^k$;
differentiate and integrate. The minimum is attained uniquely almost
surely for iid continuous normals. Condition on which coordinate is the
minimum and integrate its density against the probability that every other
coordinate lies in the interval of width $r/s$. Sum the $k$ disjoint
possibilities. If $s=0$, the range is identically zero.
\end{proof}
% Derive the closed form via the known bivariate Gaussian quadrant identity.
\begin{proposition}[Five-test closed form]\label{prop:five}
\[
a_5=\frac5{4\sqrt{\pi}}
\left(1+\frac6\pi\arcsin(1/3)\right).
\]
\end{proposition}
\begin{proof}
Integration by parts, using $\phi'=-z\phi$, gives
$a_5=20\int \phi(z)^2\Phi(z)^3\,dz$ with zero boundary term.
Since $\phi(z)^2$ is $1/(2\sqrt\pi)$ times the density of
$Y\sim N(0,1/2)$, the integral is $10/\sqrt\pi$ times
$\E\Phi(Y)^3$. For independent standard normals $U_i$ independent of $Y$,
this expectation equals the probability all three $Y-U_i$ are positive.
Their pairwise correlation is $1/3$. Expand the product of
$(1+\operatorname{sgn}(Y-U_i))/2$: the odd sign moments vanish by joint
symmetry and each pair sign expectation is $2\arcsin(1/3)/\pi$ by
the Gaussian quadrant identity \citep{tong}. Therefore the orthant
probability is $(1+6\arcsin(1/3)/\pi)/8$, proving the formula.
\end{proof}
For five tests at correlation 0.6 the exact expression evaluates to
approximately 22.0657 IQ points.
\chapter{Identification failures and sensitivity}
% Mixtures preserve marginal normality but not the joint law.
\begin{proposition}[Same marginal normals, different gaps]\label{prop:mixtures}
Let a random correlation $C$ select a conditional Gaussian
equicorrelation model with fixed $\mu,\sigma$ and $\E C=\rho$.
The unconditional marginal laws are normal, and each pair has
correlation $\rho$. The absolute gap and range means equal respectively
$2\sigma\E\sqrt{1-C}/\sqrt\pi$ and $2\sigma a_k\E\sqrt{1-C}$.
Unless $C$ is constant these are strictly smaller than the corresponding
single Gaussian answers.
\end{proposition}
\begin{proof}
The conditional marginal law does not depend on $C$. Total covariance
gives $\sigma^2\E C$ because the conditional means are constant.
Apply iterated expectation to \cref{prop:gap,prop:range}; strict Jensen
inequality for the strictly concave square root gives the comparison.
\end{proof}
A shared/independent mixture chooses $C=1$ with probability 0.6 and $0$
otherwise. A smooth mixture chooses 0.3 and 0.9 with equal probabilities.
At $k=5$ their approximate range means are 13.96 and 20.11 respectively.
% Average correlation cannot encode the entire covariance matrix.
\begin{proposition}[Average correlation is insufficient]\label{prop:block}
Take three copies of a normal score $U$ and two copies of $V$, with
common mean and SD $\sigma$ and correlation $1/3$. Their average
pairwise correlation is 0.6, but their expected range is
$2\sigma\sqrt{(2/3)/\pi}$.
\end{proposition}
\begin{proof}
Four of the ten pairs have correlation one and six have correlation
$1/3$, giving $(4+6/3)/10=0.6$. Pointwise the range is $|U-V|$.
Apply \cref{prop:gap}. The average pairwise gap instead averages
four zero gaps and six cross-group gaps.
\end{proof}
% Bounds valid without specifying a copula.
\begin{proposition}[Moment-only range bound]\label{prop:bound}
For equal means, equal variances and common correlation,
$\E R_k\le\sigma\sqrt{k(k-1)(1-\rho)}$.
\end{proposition}
\begin{proof}
Pointwise $R_k^2=\max_{i<j}(X_i-X_j)^2\le
\sum_{i<j}(X_i-X_j)^2$. Take expectations, apply
\cref{prop:moments} to every pair and Cauchy--Schwarz.
\end{proof}
% Stochastic ordering when adding observations.
\begin{proposition}[Sensitivity]\label{prop:sensitivity}
For fixed admissible correlation, $\E R_k$ is nondecreasing with
$k$ over admissible dimensions. For fixed $k$ it decreases with $\rho$,
strictly before $\rho=1$, and scales linearly in $\sigma$.
\end{proposition}
\begin{proof}
Couple all independent normals by taking an infinite iid sequence.
Adding a coordinate cannot decrease its range. The remaining claims
follow from \cref{prop:range}, since $a_k>0$ for $k\ge2$.
\end{proof}
\chapter{Conditioning versus unconditional comparison}
% Gaussian regression and folded normal formula for a nonzero mean.
\begin{proposition}[Observed first score]\label{prop:conditional}
For $|\rho|<1$, given $X_1=x$ in the Gaussian two-test model,
$D$ is normal with mean $d=(1-\rho)(x-\mu)$ and SD
$t=\sigma\sqrt{1-\rho^2}$. Its conditional absolute mean is
\[
2t\phi(d/t)+d[2\Phi(d/t)-1].
\]
\end{proposition}
\begin{proof}
The residual $X_2-\mu-\rho(X_1-\mu)$ is jointly normal with $X_1$,
has covariance zero with it and variance $t^2$. Uncorrelated jointly
normal variables are independent \citep{anderson}. Subtract the
conditional score from $x$ and integrate the resulting shifted normal
density, splitting at zero and using $\phi'=-z\phi$.
The formula defines a continuous version of the regular conditional law.
\end{proof}
At $x=145$, the predicted second score has mean 127; the signed gap
mean is 18 and the conditional absolute gap mean is about 18.70.
% Independence of the residual vector from the mean is a Gaussian result.
\begin{proposition}[Range independent of battery average]\label{prop:average}
Under the Gaussian equicorrelation model, $R_k$ and $\bar X$ are independent.
\end{proposition}
\begin{proof}
The vectors $(I-P)Z$ and $PZ$ are jointly Gaussian with zero cross
covariance; hence they are independent. The range is a measurable
function of the first and the average a measurable function of the
second by \cref{prop:domain}.
\end{proof}
% Unequal-scale Gaussian law is a direct consequence of the linear contrast.
\section{Unequal scales, means, and individual uncertainty}
\begin{proposition}[Heterogeneous scores]\label{prop:heterogeneous}
Let $(X,Y)$ be jointly normal with means $\mu_X,\mu_Y$,
positive SDs $\sigma_X,\sigma_Y$ and correlation $\rho$.
Write $\delta=\mu_X-\mu_Y$ and
$v=\sigma_X^2+\sigma_Y^2-2\rho\sigma_X\sigma_Y$.
Then $X-Y\sim N(\delta,v)$, and for $v>0$
\[
\E|X-Y|=2\sqrt v\,\phi(\delta/\sqrt v)
+\delta[2\Phi(\delta/\sqrt v)-1].
\]
At $v=0$ the absolute difference is $|\delta|$ almost surely.
\end{proposition}
\begin{proof}
Bilinearity gives the variance and linearity gives the mean;
Gaussian closure identifies the law. Split the density integral at zero,
substitute $z=(x-\delta)/\sqrt v$, and use $\phi'(z)=-z\phi(z)$.
Zero variance makes the contrast equal to its mean almost surely.
\end{proof}
Correlation one does not imply agreement: for $Y=X+c$ the absolute
discrepancy is $|c|$ although the correlation is one.
\begin{proposition}[Quantiles of the equal-mean gap]\label{prop:quantile}
When $\tau>0$ is the SD of the Gaussian contrast and $0<p<1$,
its absolute-gap $p$-quantile is $\tau\Phi^{-1}((1+p)/2)$.
\end{proposition}
\begin{proof}
The CDF in \cref{prop:gap} is continuous and strictly increasing on
the nonnegative half-line. Solve its defining equation at probability $p$.
\end{proof}
The central probability interval of a randomly selected person's
gap is therefore not a confidence interval for the population mean.

% Normal truncation uses the standard first moment of a truncated normal.
\section{Selection on a high first score}
\begin{proposition}[Threshold selection]\label{prop:selection}
In the Gaussian two-test model with $\sigma>0$, $|\rho|<1$ and finite
threshold $c$, put $z=(c-\mu)/\sigma$. Then
\[
\E[D\mid X_1>c]=(1-\rho)\sigma
\frac{\phi(z)}{1-\Phi(z)}.
\]
\end{proposition}
\begin{proof}
Iterated expectation and \cref{prop:conditional} give
$(1-\rho)\E[X_1-\mu\mid X_1>c]$.
Standardize and integrate $u\phi(u)=-\phi'(u)$ from $z$ to infinity;
divide by the positive tail probability. The boundary term vanishes.
\end{proof}
This selection result concerns signed gaps in a stipulated population,
not the expected absolute gap conditional on a single known score.

% Range second moment uses the joint density of the two extremes.
\section{Second moment of the range}
\begin{proposition}[Joint extremes]\label{prop:second}
For iid standard normals, the minimum $m_k$ and maximum $M_k$
have joint density
\[
 k(k-1)\phi(u)\phi(v)[\Phi(v)-\Phi(u)]^{k-2},\qquad u<v.
\]
Consequently, for $s=\sigma\sqrt{1-\rho}$,
\[
\E R_k^2=s^2 k(k-1)
\int_{-\infty}^\infty\int_u^\infty
(v-u)^2\phi(u)\phi(v)[\Phi(v)-\Phi(u)]^{k-2}\,dv\,du.
\]
\end{proposition}
\begin{proof}
There are $k(k-1)$ distinct assignments of the two extreme indices.
Conditional on the endpoints, each remaining independent coordinate
must fall between them; multiply these probabilities and endpoint
densities and sum the disjoint assignments. Use the pointwise range
representation in \cref{prop:range} and integrate the nonnegative
squared gap. The second moment is finite since the squared range
is bounded by the sum of squared pairwise differences
(\cref{prop:bound}).
\end{proof}
\chapter{Scope and verification}
Normality, equal scales and common correlations are assumptions, not
empirical findings. Real tests may have ceiling effects, rounding,
practice effects, heterogeneous norms or selection bias.
The numerical companion checks selected consequences by simulation
and records run-specific Monte Carlo uncertainty.
Rocq coverage is currently limited to the explicit algebraic lemmas
listed in the proof register; the probability proofs above are written
deductions using standard Gaussian, expectation and order-statistic
results, not a claim that every measure-theoretic argument is kernel checked.
\bibliographystyle{plainnat}
\bibliography{references}
\appendix
\chapter{Proof register}
\begin{tabular}{ll}
\toprule
Propositions & Foundations \\
\midrule
\ref{prop:moments}, \ref{prop:bound} & Covariance, Cauchy--Schwarz \\
\ref{prop:domain}, \ref{prop:average} & Orthogonal projections, Gaussian independence \\
\ref{prop:gap}, \ref{prop:conditional} & Gaussian density, regression \\
\ref{prop:range}, \ref{prop:five} & Order statistics, quadrant identity \\
\ref{prop:mixtures}, \ref{prop:block} & Total covariance, Jensen, counting \\
\ref{prop:sensitivity} & Pointwise monotonicity, scaling \\
\ref{prop:heterogeneous}, \ref{prop:quantile} & Gaussian density, inverse CDF \\
\ref{prop:selection}, \ref{prop:second} & Truncation, joint extremes \\
\bottomrule
\end{tabular}
\ProjectReportAIAppendix
\paragraph{This analysis specifically.} ChatGPT (GPT-6.2) prepared this draft, source integration, mathematical exposition and experimental design at Anders Hellström's request. No executed Dask run, Rocq kernel check or visual PDF review is claimed for this initial revision. The general project disclosure above must not be read as evidence that those operations were performed for this book.
\end{document}
